\documentclass[10pt,a4paper]{amsart}
\usepackage[margin=19mm]{geometry}
\usepackage{amsmath,amssymb,mathtools}
\usepackage[scr=rsfs,cal=euler]{mathalfa}
\usepackage{xfrac}
\usepackage[unicode,hidelinks,bookmarks=false]{hyperref}
\newcommand{\ie}{i.e.\ }
\newcommand{\cf}{cf.\ }
\newcommand{\resp}{respectively\ }
\newcommand{\Subsection}{\subsection}
\providecommand{\nobreakdash}{\nobreak\hbox{-}}
\setlength{\emergencystretch}{2em}
\multlinegap0pt
\usepackage{bm}
\usepackage{bbm}
\usepackage{dsfont}
\usepackage{xspace}
\usepackage{cancel}
\usepackage{mathtools}
\usepackage{amsthm}
\makeatletter
\@ifundefined{theorem}{\newtheorem{theorem}{Theorem}}{}
\@ifundefined{lemma}{\newtheorem{lemma}[theorem]{Lemma}}{}
\makeatother
\newcommand{\dif}{\mathrm{d}}
\newcommand{\intx}{\int_0^t \int_{\Omega} }
\title[Initial boundary value problems for 3-d Navier-Stokes equati]{Initial boundary value problems for 3-d Navier-Stokes equations with hyperbolic heat conduction}
\author{Yuxi Hu, Reinhard Racke}
\date{Source: arXiv:2509.22029v1 (CC BY 4.0)}
\begin{document}
\maketitle

\noindent\emph{Source and licence.} Initial boundary value problems for 3-d Navier-Stokes equations with hyperbolic heat conduction. Version arXiv:2509.22029v1 (CC BY 4.0), distributed under the Creative Commons Attribution 4.0 International licence, as recorded in the arXiv licence metadata for this exact version and shown on the abstract page https://arxiv.org/abs/2509.22029v1. Full rights notice: licenses/arxiv-2509.22029-CC-BY-4.0.txt.\par\smallskip

\noindent\textit{Editorial scope.} Counted unit: the Lemma whose source label is \texttt{le6} in the article (source0.tex lines 954--960, source label \texttt{le6}), delivered together with the article's own complete original proof of it (source0.tex lines 962--1122, one \texttt{proof} environment of 161 lines). No mathematics is written, repaired or re-derived: statement and proof are the article's own text line for line. Required context retained uncounted: 1.5 (lines 101--109). Every definition, display and label of those lines is retained unchanged. Omitted from this package and not redistributed: 1--100, 110--953, 961--961, 1123--1526. Nothing inside a retained excerpt is omitted or altered: every retained body is delivered exactly as the article's lines are written, so no omission receipt of any kind accompanies this package. Citations: the retained excerpts carry no citation command at all, so no bibliography entry of the article is transcribed and no reference is redistributed. Reference adaptation: none was needed and none was made; every reference inside the delivered unit resolves inside it, so no unresolved reference or citation is produced. Printed slips kept literal: no printed word, symbol, hypothesis or number of the counted statement or of its proof was altered. Layout and class: the article's own class is \texttt{amsart} with packages including amsfonts, amsmath, CJK, amssymb, latexsym, bm, xcolor, mathrsfs, amsthm, hyperref, ulem; the wrapper is the lane's pinned \texttt{amsart} editorial wrapper at 10pt on A4, which restates the theorem-like environments the retained text needs and adds the article's own macro definitions that the retained text needs (\texttt{\textbackslash dif}, \texttt{\textbackslash intx}), with extra wrapper packages (bm, bbm, dsfont, xspace, cancel, mathtools). The delivered numbering is the unit's own. Assets: no retained span contains a figure environment, a graphics-inclusion command, a PDF-page inclusion, a table or a TikZ picture; no image file is copied into this package, the package declares no asset dependency and no image file is redistributed with it. Licence: version arXiv:2509.22029v1 of this article is distributed under the Creative Commons Attribution 4.0 International licence, as recorded in the arXiv licence metadata for this exact version and shown on the abstract page https://arxiv.org/abs/2509.22029v1. A full rights notice is delivered at licenses/arxiv-2509.22029-CC-BY-4.0.txt. Characters: every retained excerpt is pure ASCII, so the delivered engine, XeLaTeX with its default Latin Modern text font, reports no missing character.
\begin{align}\label{1.5}
\begin{cases}
	\rho_t+(\rho u)_r+\frac{2}{r} \rho u = 0, \\
	\rho u_t+\rho u u_r+ p_r =\left( \left(\frac{4}{3} \mu(\theta)+\lambda(\theta)\right)\left(u_{r}+\frac{2}{r}u \right)\right)_r ,\\
	\rho e_\theta \theta_t + (\rho u e_\theta - \frac{2a(\theta)}{Z(\theta)} q)\theta_r + p\left(u_r+\frac{2}{r} u\right) + q_r + \frac{2}{r} q = \\
	\quad  \left(\frac{2}{\tau(\theta)}+\frac{4}{r} u\right) a(\theta) q^2 + \mu(\theta)\left( 2u_r^2+\frac{4}{r^2}u^2 - \frac{2}{3} \left(u_r+\frac{2}{r} u\right)^2\right) + \lambda(\theta)\left(u_r+\frac{2}{r} u\right)^2, \\
	\tau (\theta) \rho \left(q_t+u q_r+\frac{2}{r} u q\right) + q + \kappa(\theta) \theta_r = 0.
\end{cases}
\end{align}

 \begin{lemma}\label{le6}
 There exists a constant $C$ such that
 \begin{align}
 \int_{\Omega} \left( r^2 \rho_{tr}^2+(r^2 u_{tr}^2+2u_t^2)+r^2 \theta_{tr}^2 +\tau( r^2 q_{tr}^2+2q_t^2) \right) \dif r + \intx (r^2q_{tr}^2+2q_t^2)\dif r \dif t \nonumber\\
 +\left(\frac{4}{3} \mu+\lambda\right) \intx (r^2 u_{trr}^2+u_{tr}^2) \dif r \dif t\le C(E_0+E^\frac{1}{2}(t) \int_0^t  \mathcal D(s)\dif s+E^2(t)). \label{1.40}
 \end{align}
 \end{lemma}

 \begin{proof}
 Taking derivative with respect to $(t,r)$, one gets
 \begin{align}\label{1.41}
 \begin{cases}
 \rho_{ttr}+u \rho_{trr} +\rho u_{trr}+\frac{2}{r} \rho u_{rt}-\frac{2}{r^2} \rho u_t=H_1,\\
 \rho u_{ttr} +\rho u u_{trr}+p_{trr}-\left(\left(\frac{4}{3} \mu(\theta)+\lambda(\theta)\right)\left(u_{rr}+\frac{2}{r}u_r-\frac{2}{r^2} u\right)\right)_{tr}=H_2,\\
 \rho e_\theta \theta_{ttr}+ B(\rho, u, \theta, q) \theta_{trr}+p\left(u_r+\frac{2}{r} u\right)_{tr}+\left(q_r+\frac{2}{r} q\right)_{tr}=H_3,\\
 \tau(\theta) \rho (q_{ttr}+ u q_{trr})+q_{tr}+\kappa(\theta) \theta_{trr}=H_4,
 \end{cases}
 \end{align}
 where
 \begin{align*}
 &H_1:=-\left( \rho_t u_{rr}+2 \rho_r u_{rt}+2 u_r \rho_{tr}+u_t \rho_{rr}-\frac{2}{r^2}\rho_t u+\frac{2}{r} u\rho_{tr}+\frac{2}{r} \rho_t u_r+\frac{2}{r} \rho_r u_t\right),\\
& H_2:=-\left( \rho_t u_{tr}+\rho_r u_{tt}+u_t \rho_{tr}+(\rho u)_t u_{rr}+(\rho u)_r u_{tr}+(\rho u)_{tr} u_r\right)\\
&\qquad \qquad \qquad+\left(\lambda^\prime(\theta)\theta_r\left(u_r+\frac{2}{r}u\right)+\frac{4}{3}\mu^\prime(\theta) \theta_r\left(u_r-\frac{u}{r}\right)\right)_{tr},\\
 &H_3:=-\left((\rho e_\theta)_t \theta_{tr}+(\rho e_\theta)_r \theta_{tt}+(\rho e_\theta)_{tr}\theta_t+B_t \theta_{rr}+B_r \theta_{tr}+B_{tr}\theta_r+p_t(u_r+\frac{2}{r}u)_r\right.\\
 &\qquad \qquad \qquad\left.+p_r(u_r+\frac{2}{r} u)_t+p_{tr} \left((u_r+\frac{2}{r} u)\right)\right)+\left(\left(\frac{2}{\tau(\theta)}+\frac{4}{r} u\right) a(\theta) q^2\right)_{tr}\\
&\qquad \qquad \qquad+ \left(\mu(\theta)\left( 2u_r^2+\frac{4}{r^2}u^2-\frac{2}{3} \left(u_r+\frac{2}{r} u\right)^2\right)+\lambda(\theta)\left(u_r+\frac{2}{r} u\right)^2 \right)_{tr},\\
& H_4:=-\left( (\tau(\theta) \rho)_t q_{tr}+(\tau(\theta) \rho)_r q_{tt}+(\tau(\theta) \rho)_{tr} q_t+(\tau(\theta) \rho u)_t q_{rr}+(\tau(\theta \rho u)_r q_{tr}+(\tau(\theta) \rho u)_{tr} q_r \right.\\
& \qquad\qquad\left.+\kappa^\prime(\theta)(\theta_t \theta_{rr}+2\theta_r \theta_{tr})+\kappa^{\prime\prime} \theta_t \theta_r^2+\left( \tau(\theta) \rho \frac{2}{r} u q \right)_{tr}\right).
 \end{align*}

 Multiplying equation $\eqref{1.41}_2$ by $r^2 u_{tr}$, one gets
 \begin{align}\label{1.42}
& \int_{\Omega} \rho u_{ttr} r^2 u_{tr}\dif r  +\int_{\Omega} \rho u u_{trr} r^2 u_{tr}\dif r\nonumber\\
&+\int_{\Omega} \left(p_{trr} -\left(\left(\frac{4}{3} \mu(\theta)+\lambda(\theta)\right)\left(u_{rr}+\frac{2}{r}u_r-\frac{2}{r^2} u\right)\right)_{tr}\right) r^2 u_{tr}\dif r
 =\int_{\Omega} H_2 r^2 u_{tr}\dif r.
 \end{align}
 We estimate each term in the above equation as follows.

 First,
 \begin{align*}
 \int_{\Omega} \rho r^2 ( u_{ttr} u_{tr}+u  u_{trr} u_{tr}) \dif r =\int_{\Omega} \rho r^2 \left(  \left( \frac{1}{2} u_{tr}^2 \right)_t+u \left(\frac{1}{2} u_{tr}^2\right)_r\right) \dif r=\frac{\dif}{\dif t} \int_{\Omega} \frac{1}{2} \rho r^2 u_{tr}^2 \dif r.
 \end{align*}

 Second, using the boundary condition
 $$
 \left(p_{rt}-\left(\left(\frac{4}{3} \mu(\theta)+\lambda(\theta)\right)\left(u_{rr}+\frac{2}{r}u_r-\frac{2}{r^2} u\right)\right)_{t}\right)|_{\partial \Omega}=0,
 $$
 we have
 \begin{align*}
 \int_{\Omega} \left(p_{trr} -\left(\left(\frac{4}{3} \mu(\theta)+\lambda(\theta)\right)\left(u_{rr}+\frac{2}{r}u_r-\frac{2}{r^2} u\right)\right)_{tr}\right) r^2 u_{tr}\dif r  \\
 =-\int_{\Omega} \left(p_{tr} -\left(\left(\frac{4}{3} \mu(\theta)+\lambda(\theta)\right)\left(u_{rr}+\frac{2}{r}u_r-\frac{2}{r^2} u\right)\right)_{t}\right) (r^2 u_{tr})_r \dif r .
 \end{align*}
Using equation $\eqref{1.41}_1$, we have{\small
 \begin{align*}
 &-\int_{\Omega} p_{tr} (r^2 u_{tr})_r\dif r \\
 &=-\int_{\Omega} R\left( \theta \rho_{tr}+\rho \theta_{tr}+\theta_t \rho_r+\rho_t \theta_r\right) (r^2 u_{tr})_r \dif r\\
 &=-\int_{\Omega} R(\theta \rho_{tr}+\rho \theta_{tr}) (r^2 u_{tr})_r\dif r -\frac{\dif}{\dif t} \int_{\Omega} R(\theta_t \rho_r+\rho_t\theta_r)r^2u_{rr} \dif r+\int_{\Omega} R(\theta_t \rho_r+\rho_t\theta_r)_t r^2 u_{rr} \dif r\\
& \ge \int_{\Omega} R \theta \rho_{tr} \left(\frac{1}{\rho} r^2 \rho_{ttr}+\frac{u}{\rho} r^2 \rho_{trr}-2 u_t-\frac{1}{\rho} r^2 H_1\right) \dif r-\int_{\Omega} R\rho \theta_{tr} (r^2 u_{tr})_r  \\
&\qquad\qquad\qquad-\frac{\dif}{\dif t} \int_{\Omega} R(\theta_t \rho_r+\rho_t\theta_r)r^2u_{rr} \dif r- C E^\frac{1}{2}(t) \mathcal D(t)\\
& \ge \frac{\dif}{\dif t} \int_{\Omega} \frac{R\theta}{2\rho} r^2  \rho_{tr}^2 \dif r -\int_{\Omega} R\left( \left(\frac{\theta}{\rho} r^2\right)_t+\left( \frac{\theta}{\rho} u r^2 \right)_r \right) \frac{1}{2} \rho_{tr}^2 \dif r
 -\int_{\Omega} R\rho \theta_{tr} (r^2 u_{tr})_r \dif r \\
 &\qquad\qquad\qquad -\int_{\Omega} 2 R\theta \rho_{tr} u_t \dif r -\frac{\dif}{\dif t} \int_{\Omega} R(\theta_t \rho_r+\rho_t\theta_r)r^2u_{rr} \dif r - C E^\frac{1}{2}(t) \mathcal D(t)\\
& \ge  \frac{\dif}{\dif t} \int_{\Omega} \left(\frac{R\theta}{2\rho} r^2  \rho_{tr}^2  -R(\theta_t \rho_r+\rho_t\theta_r)r^2u_{rr}\right) \dif r   -\int_{\Omega} (R\rho \theta_{tr} (r^2 u_{tr})_r + 2 R\theta \rho_{tr} u_t )\dif r -  C E^\frac{1}{2}(t) \mathcal D(t).
 \end{align*}}
 On the other hand, using the boundary condition $u_t|_{\partial_\Omega}=0$, one has
 \begin{align*}
& \int_{\Omega}\left( \left(\frac{4}{3} \mu(\theta)+\lambda(\theta)\right)\left(u_{rr}+\frac{2}{r}u_r-\frac{2}{r^2} u\right)\right)_t (r^2 u_{tr})_r \dif r\\
  &\ge \int_{\Omega} \left(\frac{4}{3} \mu(\theta)+\lambda(\theta)\right)\left(u_{trr}+\frac{2}{r}u_{tr}-\frac{2}{r^2} u_t\right) (r^2 u_{trr}+2r u_{tr}) \dif r-CE^\frac{1}{2}(t) \mathcal D(t).\\
  &=\int_{\Omega} \left(\frac{4}{3} \mu(\theta)+\lambda(\theta)\right) \left( r^2 u_{trr}^2+4u_{tr}^2+4r u_{trr} u_{tr} -2 u_t u_{trr}-\frac{4}{r} u_t u_{tr}\right)\dif r-CE^\frac{1}{2}(t) \mathcal D(t).\\
  &\ge   \int_{\Omega} \left(\frac{4}{3} \mu(\theta)+\lambda(\theta)\right)\left(\frac{1}{5} r^2 u_{trr}^2+u_{tr}^2-\frac{2}{r^2} u_t^2\right) \dif r-C E^\frac{1}{2}(t) \mathcal D(t).
  \end{align*}
 Thus, we derive that
 \begin{align}\label{1.43}
 \frac{\dif }{\dif t} \int_{\Omega} \left(\frac{R\theta}{2\rho} r^2 \rho_{tr}^2-R(\theta_t \rho_r+\rho_t\theta_r)r^2u_{rr} +\frac{1}{2} \rho r^2 u_{tr}^2 \right) \dif r -\int_{\Omega} R \rho \theta_{tr} (r^2 u_{tr})_r\dif r \nonumber\\- \int_{\Omega} 2 R\theta \rho_{tr} u_t \dif r
 + \int_{\Omega} \left(\frac{4}{3} \mu(\theta)+\lambda(\theta)\right)\left(\frac{1}{5} r^2 u_{trr}^2+u_{tr}^2-\frac{2}{r^2} u_t^2\right) \dif r \le C E^\frac{1}{2}(t) \mathcal D(t).
 \end{align}
 Multiplying equation $\eqref{1.41}_3$ by $\frac{1}{\theta} r^2 \theta_{tr}$, and integrating the results, we obtain
 \begin{align}
 \int_{\Omega} \frac{\rho e_\theta}{\theta} \theta_{ttr} r^2 \theta_{tr} \dif r+\int_{\Omega} \frac{B}{\theta} \theta_{trr} r^2 \theta_{tr} \dif r+\int_{\Omega} R \rho\left( u_r+\frac{2}{r}u\right)_{rt} r^2 \theta_{tr} \dif r \nonumber \\
 +\int_{\Omega} \frac{1}{\theta} \left( q_r+\frac{2}{r} q\right)_{rt} r^2 \theta_{tr} \dif r =\int_{\Omega} H_3 \frac{1}{\theta} r^2 \theta_{tr} \dif r. \label{1.44}
 \end{align}
 We estimate each term in the above equation as follows.
 First,
 \begin{align*}
 \int_{\Omega} \frac{\rho e_\theta}{\theta} r^2 \theta_{ttr} \theta_{tr} \dif r  \ge \frac{\dif}{\dif t} \int_{\Omega} \frac{\rho e_\theta}{2\theta} r^2 \theta_{tr}^2 \dif r -C E^\frac{1}{2}(t) \mathcal D(t).
 \end{align*}
 Using the fact that $B|_{\partial \Omega}=0$, one has
 \begin{align*}
 \int_{\Omega} \frac{B}{\theta} \theta_{trr} r^2 \theta_{tr} \dif r =-\int_{\Omega} \left(\frac{B}{\theta} r^2\right)_r \frac{1}{2}\theta_{tr}^2 \dif r \ge -C E^\frac{1}{2}(t) \mathcal D(t).
 \end{align*}
 Second,
 \begin{align*}
 \int_{\Omega}  R \rho \left(u_r+\frac{2}{r}u\right)_{rt} r^2 \theta_{tr}\dif &r=\int_{\Omega} R \rho \left( u_{trr}+\frac{2}{r}u_{tr}-\frac{2}{r^2} u_t\right) r^2 \theta_{tr} \dif r\\
 &=\int_{\Omega} R \rho (r^2 u_{rt})_r \theta_{tr} \dif r-\int_{\Omega} 2 R \rho  u_t \theta_{tr}\dif r.
 \end{align*}
 Similarly, one get
 \begin{align*}
 \int_{\Omega} \frac{1}{\theta} \left(q_r+\frac{2}{r}q\right)_{rt} r^2 \theta_{tr}\dif r =\int_{\Omega} \frac{1}{\theta} (r^2 q_{tr})_r \theta_{tr}\dif r-\int_{\Omega} \frac{2}{\theta} q_t \theta_{tr}\dif r.
 \end{align*}
 One should pay attention to the term
 \[
  \int_{\Omega} H_3 \frac{1}{\theta} r^2 \theta_{tr} \dif r,
  \] since there will appear third-order terms. The typical  third-order term can be dealt with as follows:
  \begin{align*}
 & \int_{\Omega} \left(\left(\frac{4}{3}\mu(\theta)+\lambda(\theta)\right) u_r^2\right)_{tr} r^2 \theta_{tr} \dif r\\
&  \le C\int_{\Omega} \left(\frac{4}{3}\mu(\theta)+\lambda(\theta)\right) (u_r u_{trr}+u_{tr} u_{rr}) r^2 \theta_{tr}\dif r\\
 & \le \eta \int_{\Omega} \left(\frac{4}{3}\mu(\theta)+\lambda(\theta)\right) r^2u_{trr}^2 \dif r+ C(\eta) E^\frac{1}{2}(t) \mathcal D(t),
  \end{align*}
  where we use Young's inequality and the following estimate,
  \begin{align*}
  \| \left(\frac{4}{3}\mu(\theta)+\lambda(\theta)\right)  r^2u_{tr} u_{rr} \theta_{tr}\|_{L^1}&\le C \| \sqrt{\left(\frac{4}{3}\mu(\theta)+\lambda(\theta)\right) } u_{tr}\|_{L^\infty} \|ru_{rr}\|_{L^2} \|r\theta_{tr}\|_{L^2} \\
  &\le C E^\frac{1}{2} \mathcal D(t).
  \end{align*}
 Thus, we conclude that
 \begin{align}
 \frac{\dif}{\dif t} \int_{\Omega} \frac{\rho e_\theta}{2\theta} r^2 \theta_{tr}^2 \dif t+\int_{\Omega} R \rho (r^2 u_{tr})_r \theta_{tr} \dif r-\int_{\Omega} 2 R\rho u_t \theta_{tr} \dif r +\int_{\Omega} \frac{1}{\theta} (r^2 q_{rt})_r\theta_{tr}\dif r \nonumber\\
- \eta \int_{\Omega} \left(\frac{4}{3}\mu(\theta)+\lambda(\theta)\right) r^2u_{trr}^2 \dif r -\int_{\Omega} \frac{2}{\theta} q_t \theta_{tr} \dif r \le C E^\frac{1}{2}(t) \mathcal D(t). \label{1.45}
 \end{align}
 Multiplying equation $\eqref{1.41}_4$ by $\frac{1}{\kappa(\theta)\theta} r^2 q_{tr}$, one gets
 \begin{align}
 \int_{\Omega} \frac{\tau(\theta) }{\kappa(\theta)\theta} r^2 q_{tr} \rho (q_{ttr} + u q_{trr} )\dif r +\int_{\Omega} \frac{1}{\kappa(\theta)\theta}r^2 q_{tr}^2 \dif r+\int_{\Omega} \frac{1}{\theta} r^2 q_{tr} \theta_{trr} \dif r=\int_{\Omega} \frac{1}{\kappa(\theta)\theta} r^2 q_{tr} H_4.
 \end{align}
 We estimate again each term in the above equation separately.
First, using the mass equation, one has
 \begin{align*}
  \int_{\Omega} \frac{\tau(\theta)}{\kappa(\theta)\theta}  r^2  q_{tr} \rho( q_{ttr}+u q_{trr}) \dif r&= \int_{\Omega} \frac{\tau(\theta)}{\kappa(\theta)\theta} \rho r^2 ( (\frac{1}{2} q_{tr}^2)_t+u (\frac{1}{2}q_{tr}^2)_r) \dif r\\
  &\ge \frac{\dif}{\dif t} \int_{\Omega} \frac{\tau(\theta)}{\kappa(\theta)\theta} \rho r^2  \frac{1}{2} q_{tr}^2 \dif r -CE^\frac{1}{2}(t) \mathcal D(t).
  \end{align*}
 Second, using the fact that $\theta_{tr}|_{\partial \Omega}=0$, we have
 \begin{align*}
 \int_{\Omega} \frac{1}{\theta} r^2 q_{tr}\theta_{trr}   \dif r \ge -\int_{\Omega} \frac{1}{\theta} (r^2 q_{tr})_r  \theta_{tr} \dif r -CE^\frac{1}{2}(t) \mathcal D(t).
 \end{align*}
 Thus, we derive
 \begin{align}\label{1.47}
 \frac{\dif}{\dif t} \int_{\Omega} \frac{\tau(\theta)}{2\kappa(\theta)\theta} \rho r^2 q_{tr}^2 \dif r -\int_{\Omega} \frac{1}{\theta} (r^2 q_{tr})_r \theta_{tr} \dif r+\int_{\Omega} \frac{1}{\kappa(\theta)\theta} r^2 q_{tr}^2 \dif r \le CE^\frac{1}{2}(t) \mathcal D(t).
 \end{align}
Combining \eqref{1.43}, \eqref{1.45} and \eqref{1.47}, and choosing $\eta$ sufficiently small, we derive
\begin{align}
  &\frac{\dif}{\dif t} \int_{\Omega} \left( \frac{R\theta}{2\rho} r^2 \rho_{tr}^2 -R(\theta_t \rho_r+\rho_t\theta_r)r^2u_{rr}+\frac{1}{2}\rho r^2 u_{tr}^2+\frac{\rho e_\theta}{2\theta} r^2 \theta_{tr}^2+\frac{\tau(\theta)}{2\kappa(\theta)\theta} \rho r^2 q_{tr}^2 \right)\dif r \nonumber \\
  & -\int_{\Omega}  \left(2R u_t (\theta \rho_{tr}+\rho \theta_{tr})+\frac{2}{\theta} q_t\theta_{tr}\right)\dif r
+\int_{\Omega} \frac{1}{\kappa(\theta)\theta} r^2 q_{tr}^2 \dif r  \nonumber\\
&+ \int_{\Omega} \left(\frac{4}{3} \mu(\theta)+\lambda(\theta)\right)\left(\frac{1}{10} r^2 u_{trr}^2+u_{tr}^2-\frac{2}{r^2} u_t^2\right) \dif r \le CE^\frac{1}{2}(t) \mathcal D(t). \label{1.48}
 \end{align}
 From the momentum equation, one has
 \begin{align*}
 \rho u_{tt}+\rho u u_{tr}+p_{tr}=-(\rho_t u_t+(\rho u)_t u_r) \qquad\qquad\qquad\qquad\qquad\\
+ \left(\left(\frac{4}{3} \mu(\theta)+\lambda(\theta)\right)\left(u_{rr}+\frac{2}{r}u_r-\frac{2}{r^2} u\right)+\lambda^\prime(\theta)\theta_r\left(u_r+\frac{2}{r}u\right)+\frac{4}{3}\mu^\prime(\theta) \theta_r\left(u_r-\frac{u}{r}\right)\right)_t.
 \end{align*}
 Multiplying the above equation by $2 u_t$ yields{\small
 \begin{align}\label{1.49}
 \frac{\dif}{\dif t} \int_{\Omega} \rho u_t^2\dif r+\int_{\Omega} 2 u_t R(\rho \theta_{tr}+\theta \rho_{tr})\dif r
 +\int_{\Omega} \left(\frac{4}{3} \mu(\theta)+\lambda(\theta)\right)\left(2u_{tr}^2+\frac{2}{r^2} u_t^2 \right)\dif r \le CE^\frac{1}{2}(t) \mathcal D(t),
 \end{align}}
 where we use the fact that{\small
\begin{align*}
&\int_{\Omega} \left(\left(\frac{4}{3} \mu(\theta)+\lambda(\theta)\right)\left(u_{rr}+\frac{2}{r}u_r-\frac{2}{r^2} u\right)+\lambda^\prime(\theta)\theta_r\left(u_r+\frac{2}{r}u\right)+\frac{4}{3}\mu^\prime(\theta) \theta_r\left(u_r-\frac{u}{r}\right) \right)_t2 u_t \dif r\\
&\le -\int_{\Omega} \left(\frac{4}{3} \mu(\theta)+\lambda(\theta)\right)\left(2u_{tr}^2+\frac{2}{r^2} u_t^2 \right)\dif r+CE^\frac{1}{2}(t) \mathcal D(t).
\end{align*}}
 Similarly, taking the derivative with respect to $t$ in equation $\eqref{1.5}_4$, and multiplying the result by $\frac{2}{\kappa(\theta)\theta} q_t$, one obtains
 \begin{align}\label{1.50}
 \frac{\dif}{\dif t} \int_{\Omega} \frac{\tau(\theta)}{\kappa(\theta)\theta} \rho q_t^2 \dif r+\int_{\Omega} \frac{2}{\kappa(\theta)\theta} q_t^2\dif r+\int_{\Omega} \frac{2}{\theta} \theta_{rt}q_t \dif r \le CE^\frac{1}{2}(t) \mathcal D(t).
 \end{align}
 Therefore, by using \eqref{1.48}, \eqref{1.49} and \eqref{1.50}, one derives{\small
\begin{align}
  \frac{\dif}{\dif t} \int_{\Omega} \left(\frac{R\theta}{2\rho} r^2 \rho_{tr}^2 -R(\theta_t \rho_r+\rho_t\theta_r)r^2u_{rr} + \frac{1}{2}\rho( r^2 u_{tr}^2+2u_t^2)+\frac{\rho e_\theta}{2\theta} r^2 \theta_{tr}^2+\frac{\tau(\theta)}{2\kappa(\theta)\theta} \rho (r^2 q_{tr}^2+2q_t^2) \right)\dif r \nonumber\\
+ \int_{\Omega} \left(\frac{4}{3} \mu(\theta)+\lambda(\theta)\right)\left(\frac{1}{10} r^2 u_{trr}^2+3u_{tr}^2\right) \dif r +\int_{\Omega} \frac{1}{\kappa(\theta)\theta} (r^2 q_{tr}^2+2q_t^2) \dif r \le CE^\frac{1}{2}(t) \mathcal D(t). \label{1.51}
 \end{align}}
 Integrating this over $(0, t)$, we get the desired result. The proof of Lemm \ref{le6} is finished.
 \end{proof}

\end{document}
